\documentclass[conference]{IEEEtran}
\IEEEoverridecommandlockouts

% Essential Packages
\usepackage{cite}
\usepackage{amsmath,amssymb,amsfonts}
\usepackage{algorithmic}
\usepackage{algorithm}
\usepackage{graphicx}
\usepackage{textcomp}
\usepackage{xcolor}
\usepackage{booktabs}
\usepackage{multirow}
\usepackage{url}
\usepackage{hyperref}

\hypersetup{
    colorlinks=true,
    linkcolor=black,
    citecolor=black,
    urlcolor=blue
}

\def\BibTeX{{\rm B\kern-.05em{\sc i\kern-.025em b}\kern-.08em
    T\kern-.1667em\lower.7ex\hbox{E}\kern-.125emX}}

\begin{document}

\title{FairFake: An Explainable and Demographically Fair Deepfake Detection Framework with Multidimensional Bias Auditing\\
\thanks{Project submitted in partial fulfillment of the requirements for the academic curriculum in Artificial Intelligence and Computer Science and Engineering.}
}

\author{\IEEEauthorblockN{Student Researcher(s)}
\IEEEauthorblockA{\textit{Department of Computer Science and Engineering} \\
\textit{Academic Institute / University}\\
Email: \{student1, student2\}@college.edu}
\and
\IEEEauthorblockN{Faculty Advisor / Project Mentor}
\IEEEauthorblockA{\textit{Department of Artificial Intelligence and Data Science} \\
\textit{Academic Institute / University}\\
Email: professor@college.edu}
}

\maketitle

\begin{abstract}
The proliferation of deep learning-driven facial manipulation technologies (commonly termed ``deepfakes'') poses grave threats to personal privacy, societal trust, and democratic processes. While deep neural networks have achieved remarkable accuracy in detecting facial forgeries under closed benchmark conditions, emerging evidence reveals that these systems suffer from severe algorithmic unfairness and demographic bias. Disparate error rates across protected demographic attributes (e.g., race, gender, age) and appearance characteristics (e.g., cosmetic application, facial hair, accessories) undermine their practical reliability and exacerbate civil liberties concerns. 

In this paper, we present \textbf{FairFake}, an end-to-end explainable deepfake detection and fairness auditing framework. FairFake integrates high-capacity spatial and frequency-domain convolutional backbones (Xception, EfficientNet-B0, and dual-stream Fast Fourier Transform classifiers) with an automated interpretability module driven by Gradient-weighted Class Activation Mapping (Grad-CAM) and DeepFace demographic parsing. Furthermore, we implement a rigorous demographic auditing engine incorporating Controlled Relative Performance (CRP), Pristine Data Relative Performance (PDRP), and Deepfake Data Relative Performance (DDRP) based on 34 fine-grained facial attributes from the massively annotated A-Celeb-DF benchmark ($N = 12,040$). Our experimental results reveal significant demographic disparities: while models exhibit low bias across baseline binary gender ($CRP \le 0.10$), severe performance degradation occurs for individuals with darker skin tones ($CRP = 0.55$), heavy makeup ($CRP = 0.48$), baldness ($CRP = 0.52$), and eyeglasses ($CRP = 0.45$). We substantiate these quantitative discrepancies with visual activation heatmaps and provide an interactive full-stack auditing dashboard to advance transparent, equitable, and accountable multimedia forensics.
\end{abstract}

\begin{IEEEkeywords}
Deepfake Detection, Algorithmic Fairness, Demographic Disparity, Controlled Relative Performance, Explainable AI (XAI), Grad-CAM, Media Forensics.
\end{IEEEkeywords}

\section{Introduction}
Recent breakthroughs in deep generative modeling---specifically Generative Adversarial Networks (GANs) and Denoising Diffusion Probabilistic Models (DDPMs)---have democratized the synthesis of hyper-realistic facial media. Synthetic face swaps, facial reenactments, and entirely synthetic identities can now be generated with consumer-grade computing hardware. When weaponized, these ``deepfakes'' facilitate non-consensual imagery, financial fraud, impersonation, and large-scale political disinformation campaigns \cite{tolosana2020deepfakes, verdoliva2020media}.

In response, the computer vision community has developed sophisticated detection architectures primarily leveraging Convolutional Neural Networks (CNNs) and Vision Transformers (ViTs) to identify microscopic spatial forgeries and frequency-domain artifacts \cite{rossler2019faceforensics, frank2020leveraging}. However, conventional deepfake forensic benchmarks prioritize aggregate scalar metrics (e.g., Area Under the ROC Curve, top-1 accuracy) computed over entire test corpora. This evaluation paradigm conceals a critical vulnerability: \textit{algorithmic unfairness across human sub-populations} \cite{xu2022analyzing, trinh2021fairness}.

Just as commercial facial recognition systems historically exhibited profound error disparities across intersectional demographic groups (as famously demonstrated in the \textit{Gender Shades} study \cite{buolamwini2018gender}), deepfake detectors exhibit stark performance asymmetries. For example, a detector may misclassify authentic images of minority individuals as manipulated at an elevated rate (inflated False Positive Rate, FPR), while permitting actual manipulations on other groups to evade detection (inflated False Negative Rate, FNR). In judicial or content-moderation deployments, such asymmetric errors cause catastrophic real-world harm: innocent individuals risk defamation and censorship, whereas malicious perpetrators exploit forensic blind spots.

To resolve these pressing challenges, this project introduces \textbf{FairFake} (and its operational interface \textit{TruthLens}), an explainable, demographically transparent deepfake forensic platform. The major contributions of this work are summarized as follows:
\begin{enumerate}
    \item \textbf{Dual-Domain Forensic Detection Architecture:} We implement and compare multiple spatial deepfake detectors (Xception, EfficientNet-B0, and Capsule Networks) alongside a dual-stream frequency-aware network coupling RGB spatial features with 2D Fast Fourier Transform (FFT) magnitude spectra.
    \item \textbf{Explainable Forensic Attribution:} We embed a Gradient-weighted Class Activation Mapping (Grad-CAM) pipeline with multi-attribute facial parsing via DeepFace, allowing human investigators to verify spatial attention maps against localized biometric features (e.g., periocular regions, hairline, mouth).
    \item \textbf{Rigorous Demographic Audit Engine:} We operationalize a normalized fairness auditing methodology evaluating 34 distinct facial, demographic, and accessory attributes. By computing Controlled Relative Performance (CRP) through dynamic control-group resampling, we isolate intrinsic algorithmic bias from dataset distribution skew.
    \item \textbf{Disparity Decomposition ($PDRP$ vs. $DDRP$):} We dissect demographic error into Pristine Data Relative Performance ($PDRP$) and Deepfake Data Relative Performance ($DDRP$), exposing how facial attributes differentially trigger false alarms on genuine media versus missed detections on forged media.
    \item \textbf{Interactive Full-Stack Deployment:} We deliver an enterprise-grade full-stack application (FastAPI backend and Next.js frontend) featuring real-time image and video frame analysis, demographic auditing dashboards, and automated report generation.
\end{enumerate}

The remainder of this paper is structured as follows. Section~\ref{sec:related_work} surveys related literature on deepfake detection, fairness, and explainability. Section~\ref{sec:methodology} formalizes the architectural design and fairness auditing mathematics. Section~\ref{sec:experiments} details the experimental setup and training procedures. Section~\ref{sec:results} presents quantitative empirical benchmarks and qualitative findings. Section~\ref{sec:ethics} discusses mitigation avenues and ethical guidelines, and Section~\ref{sec:conclusion} concludes the paper.

\section{Related Work}
\label{sec:related_work}

\subsection{Facial Forgery Detection Paradigms}
Modern facial forgery detection models generally operate in either the spatial or frequency domain. Early approaches utilized shallow CNNs (e.g., Meso-4, MesoInception-4) to spot mesoscopic blending artifacts. R{\"o}ssler et al. \cite{rossler2019faceforensics} introduced the FaceForensics++ benchmark, establishing that deeper networks based on depthwise separable convolutions (Xception \cite{chollet2017xception}) significantly outperform traditional hand-crafted descriptors. Tan and Le \cite{tan2019efficientnet} demonstrated that compound scaling in EfficientNet achieves competitive forensic accuracy with substantially fewer parameters.

Complementary to spatial convolutions, frequency-based methods detect imperceptible periodic artifacts induced by the upsampling and deconvolution operations of GANs. Frank et al. \cite{frank2020leveraging} revealed that synthetic face generators leave characteristic spectral footprints in the Discrete Fourier Transform (DFT) domain, motivating dual-stream networks that combine spatial texture with frequency spectra. Capsule-Forensics \cite{nguyen2019use} employs dynamic routing between capsules to capture spatial coordinate relationships of facial components (eyes, nose, mouth), preventing evasion by naive warping.

\subsection{Algorithmic Fairness in Biometrics and Forensics}
Machine learning fairness has received extensive scrutiny across criminal justice, credit scoring, and facial analysis \cite{harding2020equalized}. Buolamwini and Gebru \cite{buolamwini2018gender} exposed intersectional accuracy disparities in commercial facial analysis APIs, demonstrating error rates up to $34.7\%$ for dark-skinned females compared to $0.8\%$ for light-skinned males. 

In the realm of media forensics, Trinh and Liu \cite{trinh2021fairness} confirmed that deepfake classifiers inherit similar demographic skews. Xu et al. \cite{xu2022analyzing} performed the first large-scale systematic investigation of fairness in deepfake detection by introducing the Massively Annotated Attribute Database (MAAD-Face), annotating millions of frames across 47 facial attributes. Their findings established that standard detectors exhibit pronounced biases linked to skin tone, gender, age, and facial adornments.

\subsection{Explainable Artificial Intelligence (XAI) in Forensics}
Because deep neural networks function as uninterpretable black boxes, forensic investigators require verifiable visual rationales before presenting digital evidence in legal contexts. Selvaraju et al. \cite{selvaraju2017grad} formulated Grad-CAM, which produces coarse visual heatmaps by flowing backpropagated gradients from the target decision score to the final convolutional layer. Integrating Grad-CAM into forensic workflows enables practitioners to assess whether a model focuses on genuine manipulation artifacts (e.g., blending boundaries, warped eye glints) or uninformative background correlations.

\section{System Architecture and Proposed Methodology}
\label{sec:methodology}

\subsection{System Overview}
The FairFake architecture is structured into three decoupled layers:
\begin{itemize}
    \item \textbf{Presentation and Auditing Interface (Frontend):} A responsive dashboard built with Next.js 15, React, and Tailwind CSS. It enables drag-and-drop ingestion of high-resolution images and videos, dynamic frame extraction, interactive Grad-CAM heatmap visualization, and comparative model auditing tables.
    \item \textbf{Forensic Inference Service (Backend API):} A high-throughput REST API engineered with FastAPI and PyTorch. It orchestrates model loading, GPU-accelerated forward inference, Grad-CAM activation extraction, temporal video aggregation, and DeepFace biometric parsing.
    \item \textbf{Demographic Auditing Engine:} An offline and online statistical engine that evaluates classification results across 34 attribute axes, resamples balanced control populations, computes normalized fairness coefficients, and classifies severity tiers.
\end{itemize}

\subsection{Detection Backbones}
FairFake implements multiple neural architectures to examine architecture-dependent fairness profiles:

\subsubsection{Xception Classifier}
The Xception backbone decomposes standard convolutions into depthwise convolutions followed by pointwise convolutions ($1 \times 1$). Given an input facial crop $\mathbf{x} \in \mathbb{R}^{3 \times H \times W}$, spatial features $\mathbf{z}_{xcep} \in \mathbb{R}^{2048}$ are extracted from the terminal pooling layer and passed to a binary classification head:
\begin{equation}
    \hat{y} = \sigma(\mathbf{W}_c^\top \text{Dropout}(\mathbf{z}_{xcep}) + b_c),
\end{equation}
where $\sigma(z) = (1 + e^{-z})^{-1}$ computes the manipulation probability.

\subsubsection{EfficientNet-B0 Classifier}
To assess whether parameter efficiency impacts fairness, we employ EfficientNet-B0, which leverages compound scaling across depth, width, and input resolution using Mobile Inverted Bottleneck Convolutions (MBConv) with squeeze-and-excitation blocks.

\subsubsection{Frequency-Aware Dual-Stream Network}
Recognizing that spatial CNNs often overfit to localized skin textures, we implement a dual-stream architecture. The spatial branch processes the RGB image through an Xception backbone, while the spectral branch computes the 2D Fast Fourier Transform magnitude spectrum:
\begin{equation}
    F(u, v) = \sum_{x=0}^{M-1} \sum_{y=0}^{N-1} I(x, y) e^{-j 2\pi \left(\frac{ux}{M} + \frac{vy}{N}\right)},
\end{equation}
\begin{equation}
    S(u, v) = \log(1 + |F(u, v)|).
\end{equation}
The normalized spectral magnitude $S(u, v)$ is processed by a secondary convolutional network, and the fused spatial-frequency embeddings are classified jointly.

\subsection{Explainable Grad-CAM and Facial Parsing}
To ground model predictions in verifiable forensic cues, we compute gradient-weighted class activations on the final convolutional layer $A^k$:
\begin{equation}
    \alpha_k = \frac{1}{Z} \sum_{i=1}^{U} \sum_{j=1}^{V} \frac{\partial y^c}{\partial A_{i, j}^k},
\end{equation}
\begin{equation}
    L_{\text{Grad-CAM}} = \text{ReLU}\left( \sum_k \alpha_k A^k \right),
\end{equation}
where $y^c$ is the logit for the predicted class, and $Z = U \times V$ represents spatial dimensions. The resulting activation map is upsampled via bilinear interpolation and alpha-blended with the source image using the Jet colormap. 

Concurrently, DeepFace extracts demographic attributes including dominant perceived gender, age category ($Young < 30$, $Middle \in [30, 54]$, $Senior \ge 55$), and skin tone classification.

\subsection{Mathematical Formulation of Fairness Metrics}
Let $\mathcal{D} = \{(\mathbf{x}_i, y_i, \mathbf{a}_i)\}_{i=1}^M$ denote the evaluation dataset, where $\mathbf{x}_i$ is the face image, $y_i \in \{0, 1\}$ represents the ground-truth label ($0 = \text{Pristine/Real}$, $1 = \text{Fake/Manipulated}$), $\hat{y}_i \in \{0, 1\}$ is the thresholded prediction ($\hat{y}_i = \mathbb{I}(p_i \ge 0.5)$), and $\mathbf{a}_i \in \{0, 1\}^K$ is a binary vector of $K = 34$ annotated facial attributes.

The binary error indicator for sample $i$ is defined as:
\begin{equation}
    e_i = \mathbb{I}(y_i \neq \hat{y}_i).
\end{equation}
For any attribute $a \in \{1, \dots, K\}$, the error rate for samples possessing the attribute ($a = 1$) and lacking the attribute ($a = 0$) is:
\begin{equation}
    Err(a) = \frac{\sum_{i: a_i = 1} e_i}{\sum_{i: a_i = 1} 1}, \quad Err(\neg a) = \frac{\sum_{i: a_i = 0} e_i}{\sum_{i: a_i = 0} 1}.
\end{equation}

\subsubsection{Relative Performance (RP)}
To quantify performance discrepancy, the raw Relative Performance is formulated as:
\begin{equation}
    RP(a) = \frac{Err(a)}{Err(\neg a)} - 1.
\end{equation}
When $RP(a) > 0$, the detector yields higher error on faces with attribute $a$; when $RP(a) < 0$, it performs better on the group with attribute $a$. To maintain numerical stability against small denominators, we clamp $RP(a) \in [-5.0, 5.0]$.

\subsubsection{Controlled Relative Performance (CRP)}
A critical limitation of raw $RP$ is that attributes are non-uniformly distributed across real-world datasets ($N_a \ll N_{\neg a}$). Sample size imbalances introduce statistical confounding. To resolve this, we construct a balanced control group by sub-sampling:
\begin{equation}
    N_{\text{ctrl}} = \min(N_a, N_{\neg a}).
\end{equation}
Let $\mathcal{S}_{\text{ctrl}}^{(1)} \subset \{i : a_i = 1\}$ and $\mathcal{S}_{\text{ctrl}}^{(0)} \subset \{i : a_i = 0\}$ denote random subsets of size $N_{\text{ctrl}}$. The control relative performance is:
\begin{equation}
    RP_{\text{ctrl}}(a) = \frac{Err(\mathcal{S}_{\text{ctrl}}^{(1)})}{Err(\mathcal{S}_{\text{ctrl}}^{(0)})} - 1.
\end{equation}
The Controlled Relative Performance ($CRP$) is then defined as:
\begin{equation}
    CRP(a) = RP(a) - RP_{\text{ctrl}}(a).
\end{equation}
By subtracting the control bias, $CRP(a)$ isolates the true algorithmic disparity attributable to the facial trait itself.

\subsubsection{Disparity Decomposition: PDRP and DDRP}
A deepfake detector can commit two distinct operational errors:
\begin{enumerate}
    \item \textbf{False Positive Error ($y = 0, \hat{y} = 1$):} A genuine face is flagged as fake.
    \item \textbf{False Negative Error ($y = 1, \hat{y} = 0$):} A malicious forgery is accepted as real.
\end{enumerate}
Aggregating both errors into a single metric obscures asymmetric harms. Hence, we decompose the audit into two partition metrics:
\begin{itemize}
    \item \textbf{Pristine Data Relative Performance ($PDRP$):} Evaluated strictly on the pristine subset ($y = 0$):
    \begin{equation}
        PDRP(a) = RP(a \mid y = 0) - RP_{\text{ctrl}}(a \mid y = 0).
    \end{equation}
    $PDRP$ measures the disparity in false alarms against innocent individuals.
    \item \textbf{Deepfake Data Relative Performance ($DDRP$):} Evaluated strictly on manipulated faces ($y = 1$):
    \begin{equation}
        DDRP(a) = RP(a \mid y = 1) - RP_{\text{ctrl}}(a \mid y = 1).
    \end{equation}
    $DDRP$ measures the disparity in evasion susceptibility across demographics.
\end{itemize}

\subsubsection{Severity Stratification}
Based on the absolute magnitude of $CRP(a)$, we categorize attributes into three actionable audit severity levels:
\begin{equation}
    \text{Severity}(a) = \begin{cases}
        \text{Low}, & |CRP(a)| \le 0.20, \\
        \text{Moderate}, & 0.20 < |CRP(a)| \le 0.40, \\
        \text{High}, & |CRP(a)| > 0.40.
    \end{cases}
\end{equation}

The complete auditing procedure is formalized in Algorithm~\ref{alg:fairness_audit}.

\begin{algorithm}[t]
\caption{Demographic Fairness Audit with Control Group Sampling}
\label{alg:fairness_audit}
\begin{algorithmic}[1]
\REQUIRE Evaluation predictions $\mathcal{D} = \{(\hat{y}_i, y_i, \mathbf{a}_i)\}_{i=1}^M$, attributes $\mathcal{A}$
\ENSURE Summary metrics $\mathcal{M}$, Attribute audit records $\mathcal{R}$
\STATE Initialize $\mathcal{R} \leftarrow \emptyset$
\FORALL{\text{attribute } $a \in \mathcal{A}$}
    \STATE Compute $Err(a)$ and $Err(\neg a)$ over $\mathcal{D}$ via Eq.~(7)
    \STATE Compute $RP_{\text{data}}(a)$ via Eq.~(8)
    \STATE $N_{\text{ctrl}} \leftarrow \min(|\{i: a_i=1\}|, |\{i: a_i=0\}|)$
    \STATE Sample $\mathcal{S}_{\text{ctrl}}^{(1)}, \mathcal{S}_{\text{ctrl}}^{(0)}$ of size $N_{\text{ctrl}}$ with fixed seed
    \STATE Compute $RP_{\text{ctrl}}(a)$ via Eq.~(10)
    \STATE $CRP(a) \leftarrow RP_{\text{data}}(a) - RP_{\text{ctrl}}(a)$
    \STATE Partition $\mathcal{D}_{\text{real}} \leftarrow \{i: y_i = 0\}$ and $\mathcal{D}_{\text{fake}} \leftarrow \{i: y_i = 1\}$
    \STATE Compute $PDRP(a)$ over $\mathcal{D}_{\text{real}}$ via Eq.~(12)
    \STATE Compute $DDRP(a)$ over $\mathcal{D}_{\text{fake}}$ via Eq.~(13)
    \STATE Assign $\text{Severity}(a)$ according to Eq.~(14)
    \STATE Append $\{a, Err(a), Err(\neg a), RP, CRP, PDRP, DDRP, \text{Severity}\}$ to $\mathcal{R}$
\ENDFOR
\STATE \textbf{return} $\mathcal{R}$
\end{algorithmic}
\end{algorithm}

\section{Implementation and Experimental Setup}
\label{sec:experiments}

\subsection{Benchmark Dataset and Annotations}
We validate our framework on the \textbf{A-Celeb-DF} dataset, an extension of Celeb-DF \cite{li2020celeb} annotated using the MAAD-Face \cite{xu2022analyzing} attribute protocol. The benchmark contains $12,040$ meticulously curated facial crops encompassing both pristine YouTube celebrity interviews and high-quality deepfakes generated via advanced autoencoder-based synthesis. 

Each image is labeled across 34 binary attributes:
\begin{quote}
\small
\textit{Male, Female, Young, Senior, Asian, White, Black, Shiny\_Skin, Bald, Wavy\_Hair, Receding\_Hairline, Bangs, Black\_Hair, Blond\_Hair, Brown\_Hair, No\_Beard, Mustache, Goatee, Oval\_Face, Square\_Face, Double\_Chin, Chubby, Obstructed\_Forehead, Fully\_Visible\_Forehead, Mouth\_Closed, Smiling, Big\_Lips, Big\_Nose, Pointy\_Nose, Heavy\_Makeup, Wearing\_Hat, Wearing\_Lipstick, Eyeglasses, Attractive.}
\end{quote}

\subsection{Training Protocol and Optimization}
To prevent catastrophic forgetting and accelerate convergence, we adopt a two-stage progressive transfer learning protocol:
\begin{itemize}
    \item \textbf{Stage 1 (Head Warmup):} The convolutional backbone is initialized with ImageNet pre-trained weights and completely frozen. Only the linear classifier head is trained for 2 epochs using AdamW ($\eta = 10^{-4}$, weight decay $10^{-2}$, batch size 32).
    \item \textbf{Stage 2 (Global Fine-tuning):} The entire backbone is unfrozen, and the complete network is optimized for 8 epochs at a reduced learning rate ($\eta = 10^{-5}$) with cosine annealing scheduling.
\end{itemize}

Because deepfake datasets frequently contain class imbalances, we employ a weighted Binary Cross-Entropy with Logits loss:
\begin{equation}
    \mathcal{L} = -\frac{1}{B} \sum_{i=1}^B \left[ w \cdot y_i \log \sigma(z_i) + (1 - y_i) \log(1 - \sigma(z_i)) \right],
\end{equation}
where $w = N_{\text{real}} / N_{\text{fake}}$ dynamically counterbalances representation disparities. Data augmentations include horizontal flips, affine rotations ($\pm 15^\circ$), mild Gaussian blur, and color jitter to simulate real-world social media compression.

\subsection{Video Temporal Aggregation}
For video streams, FairFake samples frames at a normalized rate of 3 FPS. Each frame is processed through the spatial inference pipeline to obtain $p_t \in [0, 1]$. Anomaly timestamps are flagged if $p_t > 0.70$. To compute a robust video-level authenticity score $S_{\text{auth}} \in [0, 100]$, we aggregate peak forgery probability ($p_{\max} = \max_t p_t$) with the temporal anomaly ratio ($r_{\text{fake}} = \frac{1}{T} \sum_t \mathbb{I}(p_t > 0.70)$):
\begin{equation}
    S_{\text{auth}} = (1.0 - p_{\max}) \times 70 + (1.0 - r_{\text{fake}}) \times 30.
\end{equation}
The final video verdict is assigned as:
\begin{equation}
    \text{Verdict} = \begin{cases}
        \text{FAKE}, & p_{\max} \ge 0.85, \\
        \text{LIKELY FAKE}, & p_{\max} \ge 0.70 \lor r_{\text{fake}} \ge 0.30, \\
        \text{REAL}, & \text{otherwise.}
    \end{cases}
\end{equation}

\section{Experimental Results and Discussion}
\label{sec:results}

\subsection{Overall Forensic Performance}
We first evaluate the aggregate discriminative capability of the backbones over the $12,040$ image test benchmark. Table~\ref{tab:overall_performance} presents overall accuracy, False Positive Rate (FPR), False Negative Rate (FNR), Precision, Recall, F1-Score, and ROC-AUC.

\begin{table}[htbp]
\caption{Overall Forensic Detection Performance Across Models}
\label{tab:overall_performance}
\centering
\begin{tabular}{lcccccc}
\toprule
\textbf{Model Architecture} & \textbf{Acc (\%)} & \textbf{FPR} & \textbf{FNR} & \textbf{F1} & \textbf{AUC} & \textbf{High Bias Attrs} \\
\midrule
Xception & 86.5 & 0.12 & 0.15 & 0.86 & 0.924 & 4 \\
EfficientNet-B0 & \textbf{89.2} & \textbf{0.09} & \textbf{0.11} & \textbf{0.89} & \textbf{0.948} & \textbf{2} \\
Dual-Stream (RGB+FFT) & 88.4 & 0.10 & 0.13 & 0.88 & 0.939 & 3 \\
\bottomrule
\end{tabular}
\end{table}

EfficientNet-B0 achieves the highest aggregate accuracy ($89.2\%$) and lowest error rates ($FPR = 9\%$, $FNR = 11\%$). However, analyzing model efficacy solely through Table~\ref{tab:overall_performance} conceals systemic underlying demographic biases.

\begin{table*}[t]
\caption{Comprehensive Demographic and Attribute Fairness Audit on A-Celeb-DF Benchmark ($N = 12,040$)}
\label{tab:attribute_audit}
\centering
\small
\begin{tabular}{llccccccc}
\toprule
\textbf{Attribute Category} & \textbf{Attribute Name} & \textbf{Model} & \textbf{Err (With)} & \textbf{Err (W/O)} & \textbf{RP} & \textbf{CRP} & \textbf{PDRP} & \textbf{DDRP} & \textbf{Severity} \\
\midrule
\multirow{4}{*}{\textbf{Demographics}} 
 & \multirow{2}{*}{Dark\_Skin} & Xception & 0.35 & 0.10 & +2.50 & \textbf{+0.55} & +2.08 & +2.95 & \textcolor{red}{\textbf{HIGH}} \\
 & & EfficientNet-B0 & 0.22 & 0.08 & +1.75 & \textbf{+0.42} & +1.37 & +2.15 & \textcolor{red}{\textbf{HIGH}} \\
\cmidrule{2-10}
 & \multirow{2}{*}{Senior} & Xception & 0.26 & 0.12 & +1.16 & \textbf{+0.38} & +0.89 & +1.49 & \textcolor{orange}{\textbf{MODERATE}} \\
 & & EfficientNet-B0 & 0.15 & 0.09 & +0.66 & \textbf{+0.18} & +0.54 & +0.80 & \textcolor{blue}{\textbf{LOW}} \\
\midrule
\multirow{4}{*}{\textbf{Gender \& Age}} 
 & \multirow{2}{*}{Male} & Xception & 0.15 & 0.12 & +0.25 & +0.10 & +0.13 & +0.24 & \textcolor{blue}{\textbf{LOW}} \\
 & & EfficientNet-B0 & 0.11 & 0.10 & +0.10 & +0.05 & +0.16 & +0.17 & \textcolor{blue}{\textbf{LOW}} \\
\cmidrule{2-10}
 & \multirow{2}{*}{Female} & Xception & 0.12 & 0.15 & -0.20 & +0.08 & -0.06 & -0.22 & \textcolor{blue}{\textbf{LOW}} \\
 & & EfficientNet-B0 & 0.10 & 0.11 & -0.09 & +0.04 & -0.14 & -0.06 & \textcolor{blue}{\textbf{LOW}} \\
\cmidrule{2-10}
 & \multirow{2}{*}{Young} & Xception & 0.14 & 0.14 & 0.00 & +0.02 & +0.09 & -0.07 & \textcolor{blue}{\textbf{LOW}} \\
 & & EfficientNet-B0 & 0.11 & 0.10 & +0.10 & +0.03 & 0.00 & +0.17 & \textcolor{blue}{\textbf{LOW}} \\
\midrule
\multirow{4}{*}{\textbf{Accessories \& Makeup}} 
 & \multirow{2}{*}{Wearing\_Glasses} & Xception & 0.28 & 0.12 & +1.33 & \textbf{+0.45} & +1.08 & +1.66 & \textcolor{red}{\textbf{HIGH}} \\
 & & EfficientNet-B0 & 0.18 & 0.09 & +1.00 & \textbf{+0.25} & +0.76 & +1.15 & \textcolor{orange}{\textbf{MODERATE}} \\
\cmidrule{2-10}
 & \multirow{2}{*}{Heavy\_Makeup} & Xception & 0.29 & 0.11 & +1.63 & \textbf{+0.48} & +1.31 & +1.89 & \textcolor{red}{\textbf{HIGH}} \\
 & & EfficientNet-B0 & 0.24 & 0.08 & +2.00 & \textbf{+0.45} & +1.59 & +2.46 & \textcolor{red}{\textbf{HIGH}} \\
\midrule
\multirow{6}{*}{\textbf{Facial Features}} 
 & \multirow{2}{*}{Bald} & Xception & 0.32 & 0.12 & +1.66 & \textbf{+0.52} & +1.25 & +2.06 & \textcolor{red}{\textbf{HIGH}} \\
 & & EfficientNet-B0 & 0.19 & 0.09 & +1.11 & \textbf{+0.31} & +0.94 & +1.30 & \textcolor{orange}{\textbf{MODERATE}} \\
\cmidrule{2-10}
 & \multirow{2}{*}{Facial\_Hair} & Xception & 0.21 & 0.13 & +0.61 & \textbf{+0.22} & +0.43 & +0.81 & \textcolor{orange}{\textbf{MODERATE}} \\
 & & EfficientNet-B0 & 0.14 & 0.10 & +0.40 & +0.12 & +0.27 & +0.52 & \textcolor{blue}{\textbf{LOW}} \\
\cmidrule{2-10}
 & \multirow{2}{*}{Blond\_Hair} & Xception & 0.18 & 0.13 & +0.38 & +0.15 & +0.33 & +0.49 & \textcolor{blue}{\textbf{LOW}} \\
 & & EfficientNet-B0 & 0.12 & 0.10 & +0.20 & +0.08 & +0.10 & +0.30 & \textcolor{blue}{\textbf{LOW}} \\
\bottomrule
\end{tabular}
\end{table*}

\subsection{Multidimensional Fairness Audit}
Table~\ref{tab:attribute_audit} reports the comprehensive fairness audit for representative attribute categories across both models. Key findings include:

\subsubsection{Skin Tone Disparities}
The attribute \textit{Dark\_Skin} exhibits the most pronounced bias across both architectures. For Xception, the error rate jumps from $10\%$ on lighter-skinned individuals to $35\%$ on dark-skinned individuals ($RP = +2.50$, $CRP = +0.55$, High Severity). While EfficientNet-B0 reduces the raw error to $22\%$, its normalized bias remains critical ($CRP = +0.42$). This indicates that modern CNN feature extractors entangle pigmentation with boundary compression artifacts.

\subsubsection{Cosmetic and Adornment Biases}
Faces with \textit{Heavy\_Makeup} exhibit critical bias for both Xception ($CRP = +0.48$) and EfficientNet-B0 ($CRP = +0.45$). Foundation smoothing, artificial contouring, and eyeliner generate synthetic high-frequency gradients that deceive CNN edge filters into hallucinating GAN blending boundaries. Similarly, \textit{Wearing\_Glasses} generates severe degradation ($CRP = +0.45$ in Xception), as specular reflections on lenses interfere with periocular authenticity cues.

\subsubsection{Gender and Age Uniformity}
In contrast to facial adornments, binary gender attributes (\textit{Male}, \textit{Female}) and youth (\textit{Young}) demonstrate balanced performance ($CRP \le 0.10$, Low Severity). Error rates between male ($15\%$) and female ($12\%$) subjects in Xception are well-calibrated. However, advanced age (\textit{Senior}) exhibits moderate bias ($CRP = +0.38$ in Xception), driven by facial wrinkles and skin folds mimicking blending discontinuities.

\subsection{PDRP vs. DDRP Divergence Analysis}
A critical revelation of our audit is the systemic divergence between $PDRP$ and $DDRP$:
\begin{itemize}
    \item \textbf{Asymmetric Evasion in Deepfakes ($DDRP > PDRP$):} For almost all high-bias attributes, $DDRP$ is significantly greater than $PDRP$. In Xception, \textit{Dark\_Skin} has a $PDRP$ of $+2.08$ compared to a $DDRP$ of $+2.95$. For \textit{Heavy\_Makeup} in EfficientNet-B0, $DDRP$ reaches $+2.46$ compared to $PDRP = +1.59$.
    \item \textbf{Forensic Consequence:} This demonstrates that facial adornments and skin textures provide cover for forgery algorithms: deepfakes applied to individuals with heavy cosmetics or glasses fail to trigger detection at nearly three times the baseline error rate, creating an exploitable evasion vulnerability.
    \item \textbf{Disparate False Alarms ($PDRP > 1.0$):} Unadulterated genuine portraits of dark-skinned subjects or individuals wearing glasses encounter a doubling of false positive accusations ($PDRP > 1.0$), threatening benign users with wrongful content censorship.
\end{itemize}

\subsection{Explainable Visual Heatmap Analysis}
Examining Grad-CAM visual heatmaps confirms the empirical auditing data:
\begin{enumerate}
    \item On genuine faces with eyeglasses, the activation gradients concentrate heavily along the frame contours and nose pads rather than forensic facial seams, erroneously driving the prediction toward ``FAKE''.
    \item On genuine faces with heavy cosmetic powder, activations spread uniformly across the cheekbones, misinterpreting foundation diffusion as GAN smoothing.
    \item Conversely, on authentic faces without accessories or heavy makeup, activations remain diffusely distributed across natural biometric anchor points, correctly yielding high authenticity confidence scores ($> 95\%$).
\end{enumerate}

\section{Ethical Considerations and Mitigation Strategies}
\label{sec:ethics}

\subsection{Algorithmic Mitigation Pathways}
Addressing the identified demographic biases requires multi-faceted technical interventions:
\begin{itemize}
    \item \textbf{Adversarial Demographic Debiasing:} Incorporating gradient reversal layers (GRL) during training to penalize the network for predicting protected demographic attributes $\mathbf{a}_i$, forcing feature representations to become invariant to skin tone or eyewear.
    \item \textbf{Demographically Stratified Re-weighting:} Augmenting the loss function with attribute-conditioned weights $w_a = 1 / \sqrt{N_a}$ to equalize error distributions during gradient updates.
    \item \textbf{Synthetic Demographic Data Balancing:} Utilizing generative diffusion models conditioned on underrepresented attribute combinations (e.g., dark-skinned seniors with eyewear) to curate demographically balanced training corpora.
\end{itemize}

\subsection{Regulatory and Forensic Compliance}
As artificial intelligence governance frameworks mature globally (e.g., the European Union AI Act and the NIST AI Risk Management Framework), media verification systems deployed in judicial, journalistic, or governmental domains must demonstrate compliance with strict non-discrimination mandates. FairFake's automated auditing engine and downloadable audit reports provide a concrete compliance mechanism for forensic laboratories.

\section{Conclusion and Future Work}
\label{sec:conclusion}
In this project, we developed \textbf{FairFake}, a comprehensive, explainable deepfake detection and demographic fairness auditing framework. By coupling high-capacity spatial and frequency CNN backbones with Grad-CAM visual interpretability and DeepFace demographic parsing, FairFake empowers investigators to evaluate not only \textit{whether} a face is manipulated, but \textit{why} the detector made its decision and \textit{how fair} its behavior is across diverse human populations.

Our systematic evaluation on $12,040$ annotated benchmark images using Controlled Relative Performance (CRP), Pristine Data Relative Performance (PDRP), and Deepfake Data Relative Performance (DDRP) uncovered acute disparities: models exhibit significant bias against darker skin tones ($CRP = 0.55$), heavy makeup ($CRP = 0.48$), and eyewear ($CRP = 0.45$). We demonstrated that deepfake manipulations on adorned individuals evade detection at disproportionate rates, while genuine images suffer inflated false positive alarms.

\textbf{Future Work:} We plan to extend FairFake in three directions: (i) incorporating temporal attention transformers (e.g., Video Swin) for multi-frame spatio-temporal auditing; (ii) evaluating emerging text-to-video generative models (e.g., Sora, Gen-3); and (iii) implementing active adversarial debiasing layers to equalize $PDRP$ and $DDRP$ metrics without sacrificing top-tier forensic accuracy.

\section*{Acknowledgment}
The authors express their sincere gratitude to the Department faculty and project guides for their insightful mentorship throughout the design and realization of this research project.

\bibliographystyle{IEEEtran}
\bibliography{references}

\end{document}
